\documentclass{situationist-preprint}

\title{Camus's Sisyphus Reconsidered}
\msid{EXI-04-b}
\msfamily{EXI-04}
\msdiscipline{Meaning, Death and the Absurd}
\mstype{Companion manuscript}
\mskeywords{Camus; Sisyphus; absurd; futility; revolt; irony}

\begin{document}
\maketitle

\begin{abstract}
Camus closes \emph{The Myth of Sisyphus} with the claim that one must imagine Sisyphus happy, and offers the condemned king as the hero of the absurd. EXI-04-a argues that life is absurd in Nagel's sense and that the fitting response is irony, not revolt. We reconsider the figure of Sisyphus in the light of that result. We argue that his situation is not a model of the absurd but of futility, which is a different condition; that the human situation shares his lucidity without sharing his futility; and that what is defensible in Camus's image, the persistence of serious effort without a vindication, is better described as irony than as scorn. Sisyphus is a good emblem of something, but not of the human condition.
\end{abstract}

\section{Two conditions}

Camus's Sisyphus is condemned by the gods to push a rock up a mountain from which it always rolls back. He is the absurd hero because he is lucid about his fate and does not submit to it inwardly: there is no fate that cannot be surmounted by scorn, and the struggle toward the heights is enough to fill a heart \citep{Camus1942}. The image is offered as a picture of the human situation.

EXI-04-a defines the absurd as the conjunction of two features: the subject takes some ends seriously (Definition~2.1), and from the detached standpoint (Definition~2.2) no justification of that seriousness is available that is not itself open to the same question (Definition~2.3). Futility is a different condition.

\begin{definition}\label{d:futile}
An activity is \emph{futile} if it is directed at an outcome that it cannot bring about, or at an outcome that has no value even if brought about.
\end{definition}

Sisyphus's labor is futile in both senses: the rock never stays at the top, and its staying there would serve nothing. The human situation, on Nagel's analysis, is not futile in either sense. Our projects often succeed, and many of their objects are worthwhile; what the detached standpoint discloses is not that they fail or serve nothing, but that no justification of our concern with them is immune from further questioning \citep{Nagel1971}.

\begin{proposition}\label{p:distinct}
Absurdity and futility are independent. A life can be absurd without being futile, and a task can be futile in a life that is not absurd.
\end{proposition}

\begin{proof}[Argument]
By Theorem~3.3 of EXI-04-a, the life of every reflective agent is absurd, including the life of an agent whose projects succeed and whose objects are worthwhile, such as a physician whose patients recover; that life is not futile. Conversely, an unreflective agent could be set to a futile task without ever taking up the detached standpoint, and her life would then not be absurd in the sense of Definition~2.3 of EXI-04-a, since nothing in her situation raises the question of justification.
\end{proof}

Camus's image runs the two together. Its force comes from the futility of the task, which is vivid; its claim concerns the absurdity of life, which is a different thing. The consequence is that the image overstates the bleakness of the human condition. Sisyphus's lucidity is the recognition that his task achieves nothing and is worth nothing. Ours is only the recognition that our concern with what we do cannot be vindicated from outside it.

\section{The happiness of Sisyphus}

Remark~3.7 of EXI-04-a grants Camus the lucidity of Sisyphus and rejects the framing: his scorn is directed at gods who imposed the punishment and is fitting there, but the human situation contains no such authority, and transposed to it scorn has no object. Proposition~3.5 of EXI-04-a concludes that revolt is not a fitting response to the absurd.

What then remains of the claim that Sisyphus is happy? We distinguish two readings.

On the first, Sisyphus is happy because he defies the gods. This is fitting for Sisyphus and not for us, by Remark~3.7 of EXI-04-a. An indifferent universe is not an adversary (EXI-06), and defiance of it is a response to a hostility it does not have.

On the second, Sisyphus is happy because he gives himself to the effort while knowing that nothing outside the effort vindicates it. This is the attitude EXI-04-a calls irony: a full engagement with one's ends combined with a clear perception of the discrepancy between the seriousness of that engagement and the support available for it from the most detached standpoint (Definition~2.4 of EXI-04-a). By Theorem~3.6 of EXI-04-a it is the fitting response to the absurd.

\begin{proposition}\label{p:irony}
The defensible content of Camus's conclusion is the claim that serious engagement with one's ends is compatible with lucidity about the absence of their vindication. This is the ironic response, and not revolt.
\end{proposition}

The proposition preserves what is admirable in the figure, his refusal either to despair or to deceive himself, and drops what misdescribes our situation, the presence of an adversary to be scorned.

\section{Sisyphus and meaning}

A further consequence concerns meaning. If Sisyphus's task is futile in the second sense of Definition~\ref{d:futile}, his life is not meaningful however happy he is, since meaning requires engagement with something worthwhile (EXI-17-a, Theorem~3.3). EXI-17-a makes the point against Taylor's version of the case, in which the gods implant in Sisyphus a desire to roll stones: such a Sisyphus is fulfilled but his life is not meaningful, because his object is not worthwhile \citep{Taylor1970}. Camus's happy Sisyphus is in the same position. His happiness may be real; his life, as Camus describes it, is a paradigm of the absence of meaning rather than of its presence in a godless world.

The human situation differs here as well. Our ends include worthwhile objects, and engagement with them confers meaning (EXI-01). The absurdity of our lives, on the analysis of EXI-04-a, does not touch this: the detached standpoint shows that our concern lacks an external vindication, not that its objects lack worth \citep{Wolf2010}.

\section{Objections}

\begin{objection}
Camus intended the futility of Sisyphus's task as a figure for the futility of all human effort in the face of death. Death undoes everything, as the rock rolls back.
\end{objection}

\begin{reply}
Death ends lives and many projects; it does not undo the value of what was done. A patient cured, a child raised or a truth discovered remains so after the death of the one who did it. Death limits a life without making its activities futile in either sense of Definition~\ref{d:futile}, and the figure of the rock, which literally undoes every effort, misrepresents it.
\end{reply}

\begin{objection}
Irony is a cold attitude, and Camus's revolt captures the passion with which a lucid person goes on living. Replacing revolt with irony loses that passion.
\end{objection}

\begin{reply}
Irony in the sense of EXI-04-a does not reduce engagement: it is compatible with full seriousness about one's ends, since it qualifies only the agent's view of the external support for that seriousness. The passion Camus admires belongs to the engagement and survives the change of attitude. What is lost is the antagonism, and it was misplaced.
\end{reply}

\section{Conclusion}

Sisyphus is an emblem of futility borne with lucidity, and in that role he is admirable. He is not an emblem of the human condition, which is absurd without being futile, faces no adversary, and contains objects worth the effort we give them. What Camus rightly saw in him, unbroken engagement without consolation, is irony, the response that EXI-04-a shows to be fitting.

\begin{thebibliography}{9}\small
\bibitem[Camus(1942)]{Camus1942} Camus, A. (1942). \emph{Le Mythe de Sisyphe}. Gallimard. English translation by J. O'Brien, \emph{The Myth of Sisyphus}, Hamish Hamilton, 1955.
\bibitem[Nagel(1971)]{Nagel1971} Nagel, T. (1971). The absurd. \emph{The Journal of Philosophy}, 68(20), 716--727.
\bibitem[Taylor(1970)]{Taylor1970} Taylor, R. (1970). \emph{Good and Evil: A New Direction}. Macmillan.
\bibitem[Wolf(2010)]{Wolf2010} Wolf, S. (2010). \emph{Meaning in Life and Why It Matters}. Princeton University Press.
\end{thebibliography}

\end{document}
